\section{Phụ thuộc hàm và Khóa ứng viên}

Phân tích được thực hiện trên từng quan hệ, không dùng tên thuộc tính trùng nhau ở nhiều bảng như một thuộc tính chung. Kết luận chuẩn hóa chỉ áp dụng cho tập phụ thuộc hàm nghiệp vụ được nêu dưới đây; việc dùng UUID không tự bảo đảm BCNF.

\subsection{Tập phụ thuộc hàm được xét}
Ký hiệu $R$ là toàn bộ thuộc tính của quan hệ tương ứng. Trong mỗi dòng dưới đây, từng vế trái liệt kê xác định mọi thuộc tính còn lại của $R$.
\begin{itemize}
    \item TENANT: TenantID $\rightarrow R$.
    \item MERCHANT: MerchantID $\rightarrow R$.
    \item ROLE: RoleID $\rightarrow R$.
    \item USER: UserID $\rightarrow R$; Email $\rightarrow R$, với $R$ gồm UserID, MerchantID, Email, PasswordHash, RoleID, ManagerID.
    \item CUSTOMER: CustomerID $\rightarrow R$.
    \item CAMPAIGN: CampaignID $\rightarrow R$.
    \item REWARD\_RULE: RuleID $\rightarrow R$.
    \item CREDIT: CreditID $\rightarrow R$; (CustomerID, MerchantID) $\rightarrow R$.
    \item API\_CREDENTIAL: KeyID $\rightarrow R$; ApiKey $\rightarrow R$.
    \item CREDIT\_TRANSACTION: TransactionID $\rightarrow R$, bao gồm cả CreatedAt.
\end{itemize}
Các khóa nghiệp vụ Email của USER, ApiKey và cặp (CustomerID, MerchantID) đều NOT NULL và UNIQUE. WalletAddress và SuiTxDigest chỉ duy nhất trên các dòng khác NULL; không coi chúng là khóa ứng viên của toàn quan hệ theo mô hình quan hệ cổ điển.

Không giả định RoleName, Customer.Email, CampaignName hoặc (CampaignID, ActionType) duy nhất. Không giả định MerchantID xác định RoleID hoặc ManagerID. Nếu bổ sung phụ thuộc hàm nghiệp vụ khác, cần kiểm tra chuẩn hóa lại.

\subsection{Bao đóng và tính tối thiểu}
Với USER, $\{UserID\}^{+}=R_{User}$ và $\{Email\}^{+}=R_{User}$. Không có phụ thuộc từ tập rỗng, nên cả hai tập đơn đều tối thiểu và là khóa ứng viên.

Với CREDIT, $\{CreditID\}^{+}=R_{Credit}$ và $\{CustomerID,MerchantID\}^{+}=R_{Credit}$. CustomerID riêng không xác định ví vì một khách có ví ở nhiều Merchant; MerchantID riêng không xác định ví vì một cửa hàng phục vụ nhiều khách. Do đó cặp là khóa ứng viên tối thiểu. CreditID được chọn làm PK để giao dịch tham chiếu gọn hơn; UNIQUE trên cặp vẫn bắt buộc.

Với API\_CREDENTIAL, KeyID và ApiKey là hai khóa ứng viên đơn, mỗi khóa có bao đóng bằng toàn quan hệ. Với các bảng còn lại, ID ở dòng tương ứng có bao đóng bằng toàn quan hệ và là khóa ứng viên đơn theo tập phụ thuộc hàm đã xét.

\section{Chuẩn hóa Dữ liệu}
Một quan hệ đạt BCNF nếu mọi phụ thuộc hàm không tầm thường $X \rightarrow Y$ có $X$ là siêu khóa. Trong tập phụ thuộc trên, tất cả vế trái đều là khóa ứng viên của chính quan hệ đó. Theo tập phụ thuộc này, mười quan hệ đáp ứng BCNF.

Kết luận này giảm dị thường cập nhật do các phụ thuộc hàm đã xét, không thay thế ràng buộc miền giá trị, kiểm soát đồng thời hoặc toàn vẹn liên bảng. AvailableBalance là số dư lưu sẵn để tra cứu; phải được tầng dịch vụ duy trì nhất quán với sổ cái, không thể suy ra tính nhất quán đó chỉ từ BCNF.


\subsection{Cách kiểm tra dạng chuẩn}
1NF yêu cầu mỗi ô chứa một giá trị thuộc miền đã chọn, không chứa danh sách lặp. 2NF yêu cầu đạt 1NF và thuộc tính không khóa không phụ thuộc vào một phần thực sự của khóa ứng viên. 3NF yêu cầu với mỗi phụ thuộc hàm không tầm thường, vế trái là siêu khóa hoặc thuộc tính vế phải là thuộc tính khóa. BCNF yêu cầu vế trái luôn là siêu khóa.

Ví dụ, nếu đưa TenantName vào MERCHANT thì xuất hiện MerchantID xác định TenantID và TenantID xác định TenantName. Một tổ chức có nhiều cửa hàng sẽ khiến tên tổ chức bị lưu lặp. Giữ TenantName ở TENANT và tham chiếu bằng TenantID giúp tránh cập nhật tên ở nhiều dòng MERCHANT.
